\begin{corollary}[6.12.3]
\label{IV.2.s6.12.3-fr}
Sea $X$ un preesquema localmente noetheriano. Las condiciones siguientes son
equivalentes:
\begin{enumerate}
\item[\textup{a)}] Para todo subpreesquema $Y$ de $X$,
$\operatorname{Reg}(Y)$ es abierto en $Y$.

\oldpage[IV-2]{164}

\item[\textup{b)}] Para todo subpreesquema cerrado íntegro $Y$ de $X$,
$\operatorname{Reg}(Y)$ contiene una parte abierta no vacía de $Y$.
\end{enumerate}
\end{corollary}

\begin{proof}
Está claro que \emph{a)} implica \emph{b)}, pues, si $Y$ es íntegro e $y$ es su punto
genérico, $\mathcal O_{Y,\boldsymbol{y}}$ es un cuerpo, y por tanto
$y\in\operatorname{Reg}(Y)$. Recíprocamente, para ver que \emph{b)} implica \emph{a)},
consideremos un subpreesquema cerrado íntegro $Z$ de $Y$ de punto genérico
$z$; si $Y'$ es el subpreesquema íntegro de $X$ cuyo espacio subyacente es la
clausura $Y'=\overline{\{z\}}$ de $z$ en $X$, $Z$ es abierto en $Y'$ y el
subpreesquema $Z$ de $Y'$, al ser reducido, está inducido por $Y'$ sobre el
abierto $Z$ del espacio subyacente a $Y'$. Ahora bien, la hipótesis implica
que $\operatorname{Reg}(Y')$ es una vecindad de $z$ en $Y'$, por lo que
$\operatorname{Reg}(Z)$ es una vecindad de $z$ en $Z$; basta entonces aplicar
(6.12.2) sustituyendo $X$ por $Y$ e $Y$ por $Z$.
\end{proof}
